\documentclass[11pt]{article}
\usepackage[margin=1in]{geometry}
\usepackage{lmodern,microtype}
\title{Q03 Comprehension Check}
\date{}
\begin{document}
\maketitle
\textbf{Question:} Explain in two or three sentences why ``train the model on more good responses'' (pure SFT at scale) cannot replace preference optimization --- and identify the one component of the RLHF loop that DPO eliminates entirely.

\vspace{1em}
\textbf{Answer:} SFT clones demonstrated behavior: it can only make the model more like the demonstrations, but quality is defined by \emph{comparisons between candidate responses}, which a per-sequence likelihood cannot express (there is no gradient that says ``less like $B$, more like $A$, by this margin''). Preference optimization turns that comparison into a directly optimized scalar objective. DPO eliminates the \emph{reward model} (and policy sampling during training): the Bradley--Terry preference model applied through the closed-form RLHF optimum yields a direct loss on the policy's own log-ratios.
\end{document}
